
\subsection*{1.2.}
Soit maintenant @@M0@@ un schéma quelconque et considérons le groupe @@M1@@ des classes de @@M2@@-modules localement libres (par la suite sous-entendu : de type fini). C'est une pré-@@M3@@-algèbre augmentée sur l'anneau @@M4@@ (V 3.9.1). L'application @@M5@@ permet d'identifier @@M6@@ par la suite avec un sous-anneau de @@M7@@. Si @@M8@@ est un @@M9@@-module localement libre, on désignera par @@M10@@ sa classe dans @@M11@@. On définit les opérateurs @@M12@@ (@@M13@@) par
@@M14@@
et par
@@M15@@
(V 2.2). Rappelons que dans la @@M16@@-filtration @@M17@@ de @@M18@@ on a @@M19@@ pour @@M20@@, et que pour @@M21@@, @@M22@@ est l'idéal engendré par les produits
@@M23@@
avec @@M24@@ ; @@M25@@ et @@M26@@ (@@M27@@ étant l'homomorphisme d'augmentation, cf. Exposé V 3.10). On a démontré (V 3.10) que @@M28@@ est même le sous-@@M29@@-module (donc le sous-groupe abélien, si @@M30@@ est connexe) engendré par les produits (1.2.3), et où on peut de plus supposer que les éléments @@M31@@ (@@M32@@) sont de la forme
@@M33@@
@@M34@@ étant un @@M35@@-module localement libre de rang @@M36@@.

